Khóa luận đã tiếp cận và giải quyết thành công bài toán tối ưu hóa danh mục đầu tư trong điều kiện thị trường tồn tại các rủi ro cực đoan. Thay vì dựa vào lý thuyết Markowitz với điểm yếu chí mạng là giả định phân phối chuẩn và tương quan tuyến tính, nghiên cứu đã tiên phong ứng dụng lý thuyết Copula để bóc tách và mô hình hóa chính xác cấu trúc phụ thuộc của 4 cổ phiếu công nghệ hàng đầu (AAPL, MSFT, AMZN, GOOGL).\\
Kết quả thực nghiệm chỉ ra rằng, chuỗi tỷ suất sinh lợi của các tài sản này có đặc tính đuôi nặng và mức độ đồng biến rất cao trong các trạng thái thị trường cực đoan. Hàm t-Student Copula đã chứng minh sự vượt trội khi nắm bắt thành công đặc tính này. Thông qua thuật toán mô phỏng Monte Carlo với 10.000 kịch bản, kết hợp cùng kỹ thuật cực tiểu hóa $CVaR$, nghiên cứu đã thiết lập được một cơ cấu tỷ trọng tối ưu mới. Danh mục này không chỉ duy trì được tiềm năng sinh lời mà còn thu hẹp đáng kể ngưỡng rủi ro tổn thất tối đa ($VaR_{95\%}$ giảm từ 1.56\% xuống 1.47\%), tạo ra một lớp phòng vệ vững chắc bảo vệ vốn đầu tư trước các đợt đứt gãy của thị trường.\\
Về mặt học thuật, kết quả này đóng góp một minh chứng định lượng rõ nét, phản bác lại sự phụ thuộc quá mức vào hệ số tương quan Pearson trong các giáo trình tài chính truyền thống, đồng thời khẳng định tính ứng dụng của Copula trong môi trường đầu tư thực tế. 
